% ======================================================================
% Ontologie der Kontinua -- Core 1.3.3
% Cross-K-Modul: crossk_k9_k10.tex
% Cross-Level-Relationen zwischen K9 und K10
% ======================================================================

\subsubsection{K9–K10-Querverbindung}
\label{sec:crossk-k9-k10}

Der Übergang $K_9 \rightarrow K_{10}$ markiert die Entstehung des
\textbf{metatheoretischen Kontinuums} aus dem theoretischen Kontinuum.
Während $K_9$ wissenschaftliche Theorien, Paradigmen, Modelle und logische
Rahmen enthält, führt $K_{10}$ ein:
\begin{itemize}
  \item Meta-Achsen zum Vergleich und zur Transformation ganzer Theoriesysteme,
  \item funktoriale Abbildungen zwischen Modellkategorien,
  \item Schwellen für Kreuztheorien-Konsistenz und Meta-Kohärenz,
  \item Flüsse von Interpretation, Übersetzung und Rekonstruktion,
  \item Meta-Zyklen, die die Evolution von Theorien regulieren,
  \item rekursive Strukturen, die $K_0 \dots K_{10}$ einschließlich ihrer selbst
        beschreiben.
\end{itemize}

$K_{10}$ ist die höchste repräsentationale Ebene im Core-Modell, die noch innerhalb
der humanen epistemischen Grenze bleibt.

% ----------------------------------------------------------------------
\subsubsection{1. Stufen und ihre Rollen}

\paragraph{K9 (theoretisches Kontinuum).}
$K_9$ liefert:
\begin{itemize}
  \item logische, ontologische und methodologische Achsen $A^9$,
  \item theoretische Potentiale $P^9$ (Erklärungswert, Kohärenz, Einfachheit),
  \item Schwellen $\Theta^9$ (Konsistenz, empirische Angemessenheit),
  \item Flüsse $J^9$ (Argumente, Beweise, Kritik),
  \item Paradigmenzyklen $C^9$ (Normalwissenschaft, Anomalie, Revolution).
\end{itemize}

Jedoch fehlen in $K_9$:
\begin{itemize}
  \item ein strukturierter Raum für den Vergleich von Theorien,
  \item funktoriale Abbildungen zwischen Modellstrukturen,
  \item Meta-Logik und Meta-Ontologie,
  \item rekursive Darstellung, die die Theorie der Theorien umfasst.
\end{itemize}

\paragraph{K10 (metatheoretisches Kontinuum).}
$K_{10}$ führt ein:
\begin{itemize}
  \item Achsen $A^{10}$ (Funktor-Typen, Meta-Logik, Ontologie-der-Ontologien),
  \item Potentiale $P^{10}$ (Meta-Kohärenz, Ausdrucksstärke, Universalität),
  \item Schwellen $\Theta^{10}$ (kategoriale Konsistenz, funktoriale
        Kompatibilität, meta-logische Grenzen),
  \item Flüsse $J^{10}$ (Interpretation, Übersetzung, Rekonstruktion,
        funktoriale Hebungen),
  \item Zyklen $C^{10}$ (Meta-Stabilität, Selbstkonsistenz,
        Paradigmenübergreifende Versöhnung),
  \item Maß $k(K_{10})$ der globalen Meta-Kohärenz.
\end{itemize}

Das definierende Merkmal von $K_{10}$ ist seine rekursive Kapazität:
\[
K_{10} \text{ beschreibt alle Stufen } K_0\dots K_{10}.
\]

% ----------------------------------------------------------------------
\subsubsection{2. Vererbte und neue Achsen sowie Schwellen}

\paragraph{Vererbte Struktur.}
Theorien in $K_9$ liefern Objekte; $K_{10}$ liefert Morphismen:
\[
\text{Models}(K_9) \rightarrow \text{Functors}(K_{10}).
\]

Symbolische Strukturen, Logik und Methodik aus $K_9$ werden zu Elementen der
kategorialen Architektur von $K_{10}$.

\paragraph{Neue Achsen in $K_{10}$.}
Dazu gehören:
\begin{itemize}
  \item $A_{\mathrm{functor}}$ — Typen von Modelltransformationen,
  \item $A_{\mathrm{meta}}$ — meta-logische und meta-ontologische Unterscheidungen,
  \item $A_{\mathrm{compat}}$ — Grade der Theorienkompatibilität,
  \item $A_{\mathrm{reflect}}$ — Achsen für Selbstbeschreibung.
\end{itemize}

\paragraph{Meta-Schwellen.}
$\Theta^{10}$ umfasst:
\begin{itemize}
  \item \textbf{Kohärenzschwelle} — diagrammsche Bedingungen müssen kommutieren,
  \item \textbf{Kompatibilitätsschwelle} — Funktoren müssen Struktur zwischen Theorien
        bewahren,
  \item \textbf{meta-logische Schwelle} — Vermeidung von Paradoxen bzw. Inkonsistenz,
  \item \textbf{Ausdrucksstärkeschwelle} — minimale Fähigkeit, Unterscheidungen
        über alle unteren Stufen zu repräsentieren.
\end{itemize}

Verletzung von $\Theta^9$ verhindert die Bildung von $K_{10}$.

% ----------------------------------------------------------------------
\subsubsection{3. Cross-Level-Flüsse, Zyklen und Spannungen}

\paragraph{Flüsse.}
Flüsse in $K_{10}$ erweitern theoretische Flüsse:
\[
J^{10} = (J_{\mathrm{interpret}},\ J_{\mathrm{translate}},\
         J_{\mathrm{reconstruct}},\ J_{\mathrm{lift}}),
\]
wobei:
\begin{itemize}
  \item $J_{\mathrm{interpret}}$ — Neuinterpretation einer Theorie in einen neuen Rahmen,
  \item $J_{\mathrm{translate}}$ — Abbildungen zwischen Formalismen,
  \item $J_{\mathrm{reconstruct}}$ — Wiederaufbau theoretischer Struktur,
  \item $J_{\mathrm{lift}}$ — Repräsentation einer Theorie als Teil einer höheren
        Kategorie oder eines Meta-Modells.
\end{itemize}

\paragraph{Zyklen.}
Meta-Zyklen $C^{10}$ umfassen:
\[
C_{\mathrm{meta}}:\
\text{Modell} \rightarrow \text{Interpretation} \rightarrow \text{Übersetzung}
\rightarrow \text{Verfeinerung} \rightarrow \text{Modell}.
\]

Diese Zyklen sichern die Stabilität des Meta-Kontinuums.

\paragraph{Spannung.}
Cross-Level-Spannung:
\[
T_{9,10} =
f(\Theta^{10} - P^{10},\ A^{10} - A^9,\
  J^{10} - J^9,\ \Theta^9 - \Theta^{10}).
\]
Überschreitet $T_{9,10}$ die dimensionsbezogene Schwelle, kollabiert $K_{10}$ in ein
nicht-meta-theoretisches Regime.

% ----------------------------------------------------------------------
\subsubsection{4. Geburts- und Todesbedingungen über die Stufen}

\paragraph{Geburt von $K_{10}$.}
Das metatheoretische Kontinuum entsteht, wenn:
\[
T_9 > \Theta_{9,\mathrm{dim}}
\quad\text{und}\quad
A^{10} \in M_{10} \setminus A(K_9).
\]

Erweiterung:
\[
\Omega(K_{10}) = \Omega(K_9) \cup \Delta\Omega_{\mathrm{meta}}.
\]

\paragraph{Todesausbreitung.}
Wenn theoretische Stabilität versagt ($\Theta^9$), kann $K_{10}$ nicht ausgebildet
werden.

Wenn Meta-Schwellen $\Theta^{10}$ versagen (funktoriale Inkompatibilität,
meta-logisches Paradox), kollabiert $K_{10}$, während $K_9$ bestehen bleibt.

% ----------------------------------------------------------------------
\subsubsection{5. Konzeptionelle Beispiele}

\paragraph{Theorienvergleich.}
$K_{10}$ ermöglicht strukturierte Vergleiche von Theorien über Funktoren:
\[
\mathcal{F} : \mathsf{Mod}(T_1) \rightarrow \mathsf{Mod}(T_2).
\]

\paragraph{Paradigmen interpretieren.}
Inkompatible Theorien in $K_9$ lassen sich in einem
Meta-Rahmenwerk in $K_{10}$ in Einklang bringen.

\paragraph{Vermeidung von Paradoxien.}
Verletzt ein Meta-Rahmenwerk die meta-logische Schwelle
$\Theta^{10}_{\mathrm{coh}}$, kollabiert das gesamte Meta-Kontinuum.

Damit kann der K9→K10-Übergang als begrenzte Zusammenfassung gelten als
Emergenz einer Meta-Ebene, die fähig ist, ganze theoretische Systeme zu vergleichen,
transformieren und zu rekonstruieren, einschließlich ihrer selbst.
